\chapter{Abkürzungsverzeichnis}
\label{kap:abkuerzungen}

Aufgelöst sind die Kürzel, die im Protokolltext tatsächlich vorkommen. Die Auflösung wiederholt die Erstnennung im Fließtext; die zugehörige Quelle steht im Quellenblock der Aufgabe, in der der Begriff eingeführt wird.

Nicht aufgenommen sind Kürzel, die die herangezogenen Quellen selbst nur als Kürzel oder als Eigennamen führen und nirgends auflösen. Das betrifft AMD-V, BIOS, HDD als Plattenformat von Parallels, SCSI, UEFI und VMDK; sie werden im Text deshalb als Namen gebraucht und nicht als Abkürzung erklärt.

\begin{glossar}

\begriff[sec:hosttaste]{ACPI}
Advanced Configuration and Power Interface, die Schnittstelle, über die das Betriebssystem unter anderem den Einschaltknopf gemeldet bekommt. VirtualBox bildet sie im Gast nach.

\begriff[sec:festplatte]{iSCSI}
Internet SCSI, ein Standard, der das Plattenprotokoll SCSI über das Netzwerk überträgt. VirtualBox bindet darüber Speicherserver als virtuelle Festplatte an.

\begriff[kap:grundlagen]{KVM}
Kernel-based Virtual Machine, die Virtualisierungsschicht im Linux-Kern.

\begriff[kap:grundlagen]{NAT}
Network Address Translation, der voreingestellte Netzwerkmodus von VirtualBox. Der Gast erreicht das Netz über eine Adressumsetzung, die VirtualBox zwischen VM und Wirtsystem einschiebt.

\begriff[kap:grundlagen]{PUEL}
Personal Use and Educational License, die Lizenz des VirtualBox Extension Pack.

\begriff[kap:grundlagen]{SATA}
Serial ATA, der Anschlussstandard, den VirtualBox je nach Gastbetriebssystem als voreingestellten Speichercontroller verwendet.

\begriff[kap:grundlagen]{SSE2}
Streaming SIMD Extensions 2, eine Befehlssatzerweiterung, die VirtualBox auf Wirtsystemen mit Intel-Prozessor voraussetzt.

\begriff[sec:festplatte]{VDI}
Virtual Disk Image, das eigene Plattenformat von VirtualBox.

\begriff[sec:festplatte]{VHD}
Virtual Hard Disk, das Plattenformat von Microsoft, das VirtualBox ebenfalls liest und schreibt.

\begriff[kap:grundlagen]{VM}
Virtuelle Maschine.

\begriff[sec:hypervisor]{VMM}
Virtual Machine Monitor, in der Fachliteratur gleichbedeutend mit Hypervisor.

\begriff[kap:grundlagen]{VT-x}
Intel Virtualization Technology, die Hardware-Virtualisierungserweiterung der Intel-Prozessoren.

\end{glossar}
